
On
poursuit la démonstration par récurrence sur le rang @@M0@@ de @@M1@@. Supposons le résultat prouvé pour les torseurs sous des tores déployés de rang @@M2@@. On écrit alors @@M3@@, avec @@M4@@ et @@M5@@; on pose @@M6@@ et @@M7@@, de sorte que @@M8@@. Soit @@M9@@; alors @@M10@@ est un torseur sous @@M11@@ et @@M12@@ est un torseur sur @@M13@@, de groupe @@M14@@ et d'algèbre l'image inverse de @@M15@@ sur @@M16@@. D'après ce qui précède, @@M17@@ est le quotient de @@M18@@ par le sous-groupe engendré par les @@M19@@ pour @@M20@@. Par hypothèse de récurrence, @@M21@@ est le quotient de @@M22@@ par le sous-groupe engendré par les @@M23@@ pour @@M24@@. Compte tenu de la relation
@@M25@@
@@M26@@ est donc le quotient de @@M27@@ par le sous-groupe engendré par les @@M28@@ pour @@M29@@, ou encore pour @@M30@@ parcourant un système de générateurs de @@M31@@.

\textbf{Corollaire 2.3.} \uline{Avec les notations et les hypothèses de} 2.1, \uline{supposons que} @@M32@@ \uline{soit régulier}; \uline{alors on a un isomorphisme canonique} :
@@M33@@
\uline{où} @@M34@@ \uline{est l'idéal de} @@M35@@ \uline{engendré par les} @@M36@@ \uline{pour} @@M37@@ \uline{parcourant un système de générateurs de} @@M38@@.

Si @@M39@@ est régulier, l'homomorphisme canonique @@M40@@ est un isomorphisme (IV 2.5). De même, @@M41@@ étant alors régulier, @@M42@@ est un isomorphisme. Le corollaire résulte donc immédiatement de 2.2.

